%%%%%%%%% BODY TEXT

\begin{figure*}[!t]
  \centering
  \includegraphics[width=\textwidth]{figures/fig01-dataset-sample.png}
  \caption{\textbf{Representative InfraOcc keyframe sample.}
  The sample illustrates synchronized roadside camera observations and
  voxel-level semantic occupancy annotations in the fixed roadside coordinate
  system. The left panel shows the semantic occupancy layout from a bird's-eye
  view, while the remaining panels associate multi-view roadside images with
  local 3D semantic occupancy regions. The visualization highlights the stable
  infrastructure layout and the sparse dynamic traffic participants observed
  from fixed infrastructure sensors.}
  \label{fig:dataset_sample}
\end{figure*}

\IEEEPARstart{I}{nfrastructure}-side perception is becoming an important pillar
of autonomous driving and intelligent transportation: cameras and LiDARs fixed
at intersections and road segments continuously observe long-term traffic
situations beyond the field of view of on-board sensors, providing scene-level
3D understanding for vehicles, traffic management, and safety monitoring. In
recent years, semantic occupancy prediction, which jointly estimates geometric
occupancy and semantic categories in a regular voxel space and unifies static
layout, dynamic participants, and free space within a single 3D
representation, has become an important paradigm for driving scene
understanding~\cite{OpenOccupancy,Occ3D,TPVFormer,SurroundOcc}. However, these
representations and benchmarks are almost exclusively designed for vehicle-side
ego-centric perception, where the occupancy volume is defined around a moving
ego vehicle and the viewpoint and coordinate reference change over
time~\cite{nuScenes,Occ3D,OpenOccupancy,SurroundOcc}. In contrast, roadside
cameras and LiDARs are fixed in place and repeatedly observe the same traffic
space over long periods~\cite{Rope3D,DAIRV2X,V2XReal,RCooper,TUMTrafV2X}. We
argue that this fixed viewpoint is not a mere change of camera placement, but
fundamentally alters the problem structure of 3D occupancy prediction: on the
vehicle side, occupancy prediction estimates a continuously changing
instantaneous scene frame by frame, whereas on the roadside the same space is
observed repeatedly, so the occupancy field is no longer a per-frame
instantaneous state but degenerates into a near-persistent static spatial
scaffold overlaid with sparse, short-lived dynamic traffic events.
Fixed-viewpoint roadside occupancy should therefore be treated as a distinct
problem with its own structure, rather than a viewpoint transfer of
vehicle-side occupancy.

Existing research, however, advances along two main lines, neither of which
addresses this structure. First, mainstream semantic occupancy benchmarks and
methods are ego-centric: whether they build BEV/voxel features through view
transformation and depth lifting~\cite{BEVDet,BEVDepth}, aggregate multi-view
and temporal context with transformer queries~\cite{BEVFormer,TPVFormer}, or
improve voxel reasoning and sparse/compact
representations~\cite{VoxFormer,OccFormer,SurroundOcc,SparseOcc,GaussianFormer},
their voxel space is defined around the ego vehicle with a moving coordinate
frame, and thus neither characterizes the long-term spatial structure under a
fixed roadside viewpoint nor provides data and diagnostic metrics that measure
static-dynamic organization. Second, existing infrastructure-side and
cooperative perception studies are largely object-centric: roadside and
cooperative perception
datasets~\cite{Rope3D,DAIRV2X,V2XReal,RCooper,TUMTrafV2X,OPV2V,V2XRadar,
zheng2026omnihd} and
methods targeting roadside-camera geometry modeling, ground priors, scenario
generalization, calibration robustness, and cooperative feature
fusion~\cite{BEVHeight,BEVHeightPP,MonoGAE,SGV3D,SIFormer,V2XViT,Where2comm}
mostly focus on 3D detection, tracking, or cooperative detection, yielding
sparse object boxes rather than a dense scene description that jointly
represents static infrastructure, dynamic participants, and free space. A few
cooperative semantic occupancy works have begun to connect V2X perception with
dense scene completion~\cite{CollaborativeSemanticOcc,SyntheticCollaborativeOcc},
but they mainly target connected-vehicle cooperation or synthetic settings, and
dense semantic occupancy under real fixed roadside sensors remains unexplored.
As a result, the core structure of ``persistent static versus transient
dynamic'' under a fixed viewpoint has neither a benchmark to measure it nor a
method to exploit it.

To fill this gap, we build \textbf{InfraOcc}, to our knowledge, the first real-world
infrastructure-side semantic occupancy benchmark. Constructing roadside
occupancy annotations is not a simple conversion of existing 3D annotations into
voxel labels: fixed roadside LiDAR is limited by a single viewpoint and
occlusions, making it hard to fully cover persistent static background such as
roads, buildings, vegetation, and guardrails, while dynamic objects in a single
roadside scan are often sparse and boundary-incomplete. InfraOcc therefore
adopts a static-dynamic decoupled annotation pipeline: a dynamic branch
aggregates dynamic geometry from roadside LiDAR sequences and temporally
consistent object tracklets, a static branch completes the static background
using the moving viewpoints of vehicle-side LiDAR, and the two are then
recomposed in a fixed roadside coordinate system with visibility-aware labeling.
We emphasize that vehicle-side LiDAR is used only for annotation construction;
benchmark inputs are always restricted to roadside cameras, roadside LiDAR, or
their combination. InfraOcc further provides a unified camera-only, LiDAR-only,
and multi-modal evaluation protocol, together with diagnostic metrics that
separate static and dynamic occupancy.

Statistical analysis of InfraOcc provides the first quantitative
characterization of this structural disparity. Temporally, static
infrastructure is highly persistent across frames (a median temporal occupancy
ratio close to 100\%), whereas dynamic participants have a median temporal
occupancy ratio of only 1.4\% per location and a 95th percentile of about
15.7\%; in terms of semantic frequency, static infrastructure accounts for
97.3\% of occupied voxels while dynamic participants account for only 2.7\%.
This shows that the fundamental difficulty of roadside occupancy is not merely
semantic long-tailedness, but the structural asymmetry between persistent static
structure and transient dynamic evidence. Under this structure, flat one-shot
voxel classification biases the model toward frequent and spatially dominant
static patterns, weakening sparse dynamic cues during voxel feature construction
and 3D context aggregation; simple loss
reweighting~\cite{FocalLoss,ClassBalancedLoss} can alleviate class imbalance but
cannot change the way static and dynamic evidence are homogeneously organized.

Building on this, we propose \textbf{ProSD-Occ}, a progressive
static-to-dynamic reasoning framework that reformulates fixed-viewpoint
occupancy prediction from ``one-shot classification of all voxels'' into
``step-by-step reasoning ordered by evidence reliability.'' Its core idea
directly mirrors the structural asymmetry above: since a roadside scene is the
superposition of a persistent static scaffold and sparse, transient dynamic
events on top of it, the model should first explain the stable, reusable part
and then infer the part that this explanation cannot account for, rather than
letting sparse dynamic cues be overwhelmed by spatially dominant static patterns
in a single classification. Concretely, ProSD-Occ first establishes a
sample-adaptive explanation of the persistent static layout, then uses this
explanation as a reference to explicitly expose residual evidence that the
static layout cannot sufficiently account for---often corresponding to sparse
yet safety-critical dynamic objects---and finally recomposes static, dynamic,
and free-space evidence into a unified semantic occupancy field. This
progressive process is realized by four interconnected modules---a Static Layout
Reasoner, Static-guided Residual Modulation, a Residual Dynamic Predictor, and
Semantic Recomposition---responsible for layout explanation, evidence exposure,
dynamic inference, and unified recomposition, respectively. In this way,
ProSD-Occ turns the static-dynamic asymmetry of the fixed viewpoint from a bias
that hampers prediction into an explicit, learnable reasoning structure: it
retains the contextual value of stable static layout while preventing it from
dominating the prediction of sparse dynamic objects.

We conduct a systematic evaluation of representative occupancy baselines on
InfraOcc. The experiments cover camera-only, LiDAR-only, and multi-modal
tracks, and report overall semantic occupancy, dynamic occupancy, static
occupancy, and geometric occupancy metrics. ProSD-Occ consistently outperforms
representative baselines under all three sensing settings, achieving 60.08,
63.66, and 65.87 mIoU$_{\mathrm{all}}$, and 28.97, 33.87, and 37.92
mIoU$_{\mathrm{dyn}}$, respectively. These results demonstrate that
static-to-dynamic evidence organization provides stable benefits for
fixed-view infrastructure occupancy prediction.
The main contributions of this paper are summarized as follows:
\begin{itemize}
\item We introduce \textbf{InfraOcc}, a real-world infrastructure-side
semantic occupancy benchmark that constructs dense voxel-level annotations in
a fixed roadside coordinate system and supports unified camera-only,
LiDAR-only, and multi-modal evaluation. 

\item Based on InfraOcc, we further provide
a systematic diagnosis of \textbf{static-dynamic asymmetry} in roadside
occupancy, revealing that the structural asymmetry between persistent static
infrastructure and sparse dynamic traffic participants is a core challenge of
this task.

\item We propose \textbf{ProSD-Occ}, a progressive static-to-dynamic
reasoning framework that explicitly converts static-dynamic asymmetry into a
learnable reasoning structure through Static Layout Reasoner,
Static-guided Residual Modulation, Residual Dynamic Predictor, and Semantic
Recomposition.

\item We compare representative occupancy baselines on InfraOcc and show that
ProSD-Occ achieves better mIoU$_{\mathrm{all}}$ and
mIoU$_{\mathrm{dyn}}$ across camera-only, LiDAR-only, and multi-modal
settings.
\end{itemize}
